\documentclass[11pt]{article}
\usepackage[margin=1in]{geometry}
\usepackage{amsmath,amssymb,amsthm}
\usepackage{hyperref}
\usepackage{booktabs}
\hypersetup{hidelinks}

\newtheorem{theorem}{Theorem}
\newtheorem{lemma}[theorem]{Lemma}
\newtheorem{corollary}[theorem]{Corollary}
\newtheorem{definition}{Definition}

\setlength{\parskip}{4pt plus 1pt minus 1pt}
\setlength{\parindent}{0pt}

\title{Uncertainty Bounding as the Basis of Intelligence:\\A Formal Theory of Bounded-Context Intelligence}
\author{Patrick McCarthy}
\date{}

\begin{document}
\maketitle

\begin{abstract}
We propose that intelligence is the \emph{driven capacity to improve}: agency $\mathcal{A}$ exercised through a higher-order function space $M$ that refines both the knowledge graph $K$ and the informed action space $F$. Formally: $I(E) = \mathcal{A} \cdot \eta_M$. Any entity persisting in a world state space $W$ operates under a physically bounded active context $C_n$. This constraint has structural consequences: intelligence cannot grow by filling $C_n$ --- it must grow by encoding proven knowledge at lower levels, freeing $C_n$ for genuine novelty. We derive from a single generative constraint --- bounded $C_n$ under survival pressure in a world with positive entropy --- that (i) knowledge is necessarily graph-structured (sub-linear retrieval requires stored adjacency), (ii) bottom-up escalation dominates top-down allocation as capability grows, (iii) the encoding hierarchy follows from cost structure rather than requiring separate explanation, (iv) retention criteria are asymmetric (uncertain benefit suffices; certain zero benefit is required for pruning), and (v) $\mathcal{A} > 0$ and $\eta_M > 0$ are the unique structural invariants of persistence. Prior theories --- FEP, AIXI, bounded rationality --- each take their constraint as given. The contribution is that these structural consequences are all derivable from the same constraint. The extended version develops proofs in full; a companion paper extends to collective intelligence and sentience.
\end{abstract}

\section{Introduction}

This framework emerged from a concrete engineering problem: generating repeatable, transposable tasks across a diverse codebase ecosystem. The formal theory describes what a self-optimizing agentic system converged to under real performance constraints. The biological and civilizational correspondences are predictions the framework makes --- not the premises it rests on.

As the world $W$ expanded, the context window became the central constraint. Encoding what you know together with how to apply it does not scale: shared knowledge cannot be maintained consistently when it is duplicated at every point of use. The knowledge graph solved this by separating \emph{why you know something} from \emph{how to apply it}: $O(N \times M) \to O(1)$ retrieval via local graph traversal.

\textbf{Central claim.} Intelligence is the driven capacity to improve --- agency $\mathcal{A}$ exercised through $M$ that continuously refines $(F, K)$:
\[I(E) = \mathcal{A} \cdot \eta_M\]
Three components are jointly necessary: \emph{drive} ($\mathcal{A}$), \emph{mechanism} ($\eta_M$), and \emph{outcome} (uncertainty reduction). $\mathcal{A} > 0$ and $\eta_M > 0$ are the only universal invariants (Theorem~\ref{thm:necessity}). Every other property of $E$ can differ arbitrarily.

\textbf{Contribution.} Prior theories identify individual aspects: FEP derives uncertainty minimization from self-organization~[Friston2010]; Hutter defines optimal behavior over computable environments~[Hutter2005]; Chollet measures skill-acquisition efficiency~[Chollet2019]; Simon establishes bounded rationality~[Simon1955]. Each takes its constraint as given. This paper identifies the \emph{generative constraint} --- bounded $C_n$ under survival pressure --- and derives the structural consequences as a connected chain: graph-structured $K$ (Theorem~\ref{thm:graph}), escalation dominance (Theorem~\ref{thm:topdown}), encoding permanence (Theorem~\ref{thm:encoding}), retention criteria (Theorem~\ref{thm:retention}), and the necessity of $\mathcal{A}$ and $\eta_M$ (Theorem~\ref{thm:necessity}). The individual results formalize known intuitions; the contribution is that they are all consequences of the same constraint.

\section{Formal Framework}

\begin{definition}[World, Entity, Uncertainty]
Let $W$ be a measurable space of world states with positive entropy. An entity $E = (K, F, \sigma, \mathcal{A})$ where $\sigma : W \to W_{obs}$ is a sensing function. The uncertainty of $E$ about $w$ is $U(w, K) = -\log P(w \mid K)$ where $P(w \mid K) \propto P_0(w) \cdot \prod_{p \in K} \ell_p(w)$ is the Bayesian posterior given knowledge graph $K$~[Shannon1948, Pearl1988]. Expected uncertainty is $H(W \mid K) = \mathbb{E}_w[U(w,K)]$.
\end{definition}

\begin{definition}[Survival Threshold]
For entity $E$ in domain $d \subseteq W$, there exists $U_{lethal}^d$ such that $U(w,K) > U_{lethal}^d \Rightarrow E$ does not survive in $d$. This threshold is imposed by $W$, not chosen.
\end{definition}

\begin{definition}[Knowledge Graph]
$K$ is a directed graph where nodes are propositions about $W$ and edges represent inferential or causal relationships~[Pearl2009]. $K$ has bounded capacity $K_{max}$. In the \emph{growth regime} ($|K| < K_{max}$), propositions accumulate; in the \emph{curation regime} ($|K| = K_{max}$), adding requires removing.
\end{definition}

\begin{definition}[Informed Action Space]
$F \subseteq (K \times W_{obs} \to \Delta A)$ is the set of informed actions: actions grounded in $K$, validated against $W$. Feedback from executing $f \in F$ at $w$ is $\phi(f,w) \in \Phi$, yielding updated $K_f = K \cup \{p_f\}$. The information gain is $G(f,K) = I(f; W \mid K) \geq 0$ by the data processing inequality~[Cover2006].
\end{definition}

\begin{definition}[Agency and Intelligence]
Agency $\mathcal{A} \in [0,1]$ is the probability that $E$ engages its learning mechanism $M$ in response to a gradient signal from $W$: the response gate. $\mathcal{A}(E) \propto \lim_{U_{lethal}^d \to \infty} \frac{dU}{dt}(E)$: an entity with $\mathcal{A} > 0$ continues learning when survival pressure is removed. The higher-order function space $M : (F,K) \times \mathcal{L} \to (F,K)$ operates on the joint $(F,K)$ pair, where $\mathcal{L}$ is a learning signal~[Thrun1998]. $M$ efficiency is $\eta_M(E) = \mathbb{E}_\mathcal{L}\!\left[\frac{\max_{f' \in F'} G(f',K') - \max_{f \in F} G(f,K)}{|\mathcal{L}|}\right]$. Intelligence: $I(E) = \mathcal{A} \cdot \eta_M$.
\end{definition}

\begin{definition}[Encoding Hierarchy and Context Capacity]\label{def:encoding}
The encoding hierarchy is a continuous space $\mathcal{E}(c,r)$ parameterized by runtime cost $c \geq 0$ and reversibility $r \in [0,1]$~[Baars1988, Kahneman2011]. Interfaces $L_0, \ldots, L_n$ mark qualitative boundaries. $L_n$ (active reasoning) has bounded capacity $C_n$. Inference quality: $\rho(L_n, f) = |K^{[f]}| / |\text{ctx}(L_n, t)|$, where $K^{[f]}$ is the relevant subgraph. Under top-down allocation with $k(t)$ delegated subproblems and per-delegation overhead $\alpha$: $\rho(L_n, f_{\text{novel}}) = |K^{[f]}| / (|K^{[f]}| + k(t) \cdot \alpha)$.
\end{definition}

\section{Knowledge-Action Inseparability}

\begin{lemma}[Proposition Absence Lethality]
For $p \in K$ grounding $f \in F$ in domain $d$: $H(W \mid K) \leq H(W \mid K \setminus \{p\})$. If $p$ is the sole grounding for the critical action in $d$, removal may cause $U > U_{lethal}^d$.
\end{lemma}

\begin{theorem}[Retention Criterion]\label{thm:retention}
The retention criterion is asymmetric: $p$ is retained when survival benefit is uncertain; pruned only when benefit is certainly zero across all reachable trajectories. In the curation regime, $p$ is displaced when a candidate $q$ satisfies $\mathbb{E}[G(f_q, K)] > \mathbb{E}[G(f_p, K)]$ at the margin.
\end{theorem}

\emph{Proof sketch.} The option value of retaining $p$ is $OV(p) = P(\exists\,\tau^* : B(p,\tau^*) > 0) \cdot \mathbb{E}[B(p,\tau^*) \mid B > 0]$. The cost of absent $p$ when needed is bounded below by entity non-survival (Lemma 1), which dominates finite maintenance costs. Retention holds whenever $P(\exists\,\tau^* : B(p,\tau^*) > 0) > 0$. Pruning requires this probability to equal zero. The asymmetry is strict. Full proof in extended version. $\square$

\textbf{Corollary} (Indexed Knowledge). $K$ is constitutively indexed by $F$: different $F$ produces different $K$ from identical observations. This formalizes Gibson's affordances~[Gibson1979]; the retention threshold extends the information bottleneck~[Tishby2011] by deriving it from survival pressure. Anderson and Schooler~[AndersonSchooler1991] confirm empirically that human memory retention mirrors the statistical structure of the environment --- precisely this prediction.

\section{Escalation and Graph Necessity}

\begin{theorem}[Escalation]\label{thm:escalation}
$f \in F$ is handled at the lowest $L_i$ such that $U(w, K^{[f]}) \approx 0$. Escalation to $L_{i+1}$ occurs iff $L_i$ fails to bound uncertainty. (Proof: monotone cost $C_{i+1} > C_i$ makes over-escalation wasteful; selection removes it.)
\end{theorem}

\begin{theorem}[Top-Down Allocation Degrades]\label{thm:topdown}
Under top-down allocation with $k(t)$ delegated subproblems:
$\rho(L_n, f_\text{novel}) = |K^{[f]}| / (|K^{[f]}| + k(t) \cdot \alpha)$.
As $|F|$ grows, $k(t) \to \infty$ and $\rho \to 0$. Recovery requires growing $C_n$ linearly with capability.
\end{theorem}

\emph{Proof.} Each delegation occupies $\alpha > 0$ context slots for specification, expected return, and integration context. Novel problem $f$ requires $|K^{[f]}|$ slots. Total: $|K^{[f]}| + k(t) \cdot \alpha$. As $k(t)$ grows, delegation overhead dominates: $\rho \to |K^{[f]}| / (k(t) \cdot \alpha) \to 0$. Growing $C_n$ is the only recovery, but larger $C_n$ admits more variance, creating a feedback loop. Under escalation (Theorem~\ref{thm:escalation}), routine problems are handled below $L_n$ with zero delegation overhead: $\rho = 1$ for novel problems. Escalation is $O(1)$ in steady state; top-down is $O(|F|/C_n)$. $\square$

This is a critique of FEP's architectural \emph{prescription}~[Friston2010, Clark2015], not its minimization objective. Precision weighting may be optimal when $C_n$ is not binding; under bounded context, escalation dominates.

\begin{theorem}[Graph Necessity]\label{thm:graph}
Graph-structured $K$ is the necessary representation under bounded $C_n$.
\end{theorem}

\emph{Proof.} Three simultaneous constraints: retrieval cost $O(|K^{[f]}|)$ (proportional to relevant subgraph, not $|K|$), context purity $\rho = 1$ (no extraneous propositions), and deduplication (shared propositions stored once). When $|K|$ reaches $C_n$, making space admits three options:

\textbf{(A) Discard.} Without relational structure, the entity cannot determine which propositions share dependencies. Discarding loses information and potentially breaks dependents --- a second-order loss.

\textbf{(B) Factor.} Extract shared structure as a named node with typed edges. Preserves information; creates $O(1)$ update propagation. Factoring \emph{is} graph construction.

\textbf{(C) Lossy compression.} Replace propositions with a summary. If dependents exist, compression without knowing adjacency silently invalidates them. With adjacency, compression can be scoped correctly --- but stored adjacency is a graph. Correct compression reduces to (B); blind compression reduces to (A).

Any representation preserving information as $|K|$ grows beyond $C_n$ must factor. Every factoring operation creates nodes with typed edges. $\square$

An independent proof from retention dynamics (Theorem 2c in the extended version) yields the same conclusion via a lower bound: without stored adjacency, dependency queries require $\Omega(n)$; cascade removal amplifies to $\Omega(n \cdot L)$. Sub-linear cost on all four retention operations (addition, retrieval, dependency query, selective removal) requires a directed graph.

\textbf{Corollary} (Selection Produces Graphs). Under selection pressure, entities whose $K$ has graph structure outperform those without in retrieval, update, and cascade. Selection produces graph-structured $K$.

\section{Gradient Descent and Encoding Permanence}

\begin{theorem}[Gradient Descent on Uncertainty]\label{thm:gradient}
Every executed $f \in F$ generates feedback $\delta = \phi(f,w) - \mathbb{E}[\phi(f,\cdot) \mid K]$: a subgradient of $U$ as a functional on $F$. $\mathcal{A}$ is the response gate: the probability $M$ engages on $\delta$.
\end{theorem}

\emph{Convergence.} $U(w,K) = H(w \mid K) \geq 0$ is a Lyapunov function. Each $M$-update with $G > 0$ strictly reduces $\mathbb{E}[U]$ (data processing inequality: $H(W \mid K') \leq H(W \mid K)$~[Cover2006]). The sequence $\{\mathbb{E}[U(w, K(t))]\}$ is non-increasing and bounded below; by monotone convergence, $\mathbb{E}[U]$ converges. Cross-grounding of new propositions with existing $K$ accelerates convergence --- adding structure can improve resolution of existing knowledge, analogous to re-normalizing a matrix. The convergence operates on scalar $U$; dimensionality of $W$ or $F$ is irrelevant.

\emph{Non-termination.} Each executed $f$ produces feedback grounding a new proposition, extending $\mathcal{T}_{reachable}$. The frontier always contains states where $H(w \mid K) > 0$; the gradient never terminates within the entity's reachable domain.

\begin{theorem}[Encoding Permanence]\label{thm:encoding}
The certainty threshold for promotion to $L_i$ is $\theta_i = 1 - \varepsilon_\text{acceptable}/C_i$, where $C_i$ is the cost of failure at $L_i$.
\end{theorem}

\emph{Proof.} Expected encoding error cost: $(1-\theta) \cdot C_i$. Promotion is safe when $(1-\theta) \cdot C_i \leq \varepsilon_\text{acceptable}$, giving $\theta_i = 1 - \varepsilon_\text{acceptable}/C_i$. Higher permanence requires higher certainty. $\square$

\textbf{Corollary} (Bounded Adaptation). Certain informed actions escape $L_n$ permanently, freeing it for the next frontier.

\section{Structural Invariants}

\begin{theorem}[Type Irreducibility of $F$ and $M$]\label{thm:type}
$F$ and $M$ are type-irreducible: $M$ cannot be a member of $F$ or vice versa. $M : (F,K) \times \mathcal{L} \to (F,K)$ operates \emph{on} $(F,K)$; if $M \in F$, then $M$ would need to take itself as input, requiring $M : (F \cup \{M\}, K) \times \mathcal{L} \to (F \cup \{M\}, K)$ --- an infinite regress in the type signature. $\square$
\end{theorem}

\begin{theorem}[Necessity of $\mathcal{A}$ and $\eta_M$]\label{thm:necessity}
For any entity $E$ persisting in $W$ under survival pressure with bounded $C_n$, in a world with positive entropy: (1) $\mathcal{A}(E) > 0$ necessarily; (2) $\eta_M(E) > 0$ necessarily; (3) every other scalar property of $E$ can equal zero for some persistent entity; (4) $I(E) = \mathcal{A} \cdot \eta_M$ is the natural measure of persistence.
\end{theorem}

\emph{Proof.} \textbf{Part 1.} Suppose $\mathcal{A} = 0$: $(F,K)$ is fixed for all $t$. $W$ has positive entropy and generates novel configurations no finite $K(0)$ anticipates. The selection argument operates on lineages: $E$ reproduces copies with the same fixed $(F(0), K(0))$; copies disperse across $W$; some encounter states where $U > U_{lethal}^d$. No copy can adapt. Niches shift as $W$ evolves. Selection removes the $\mathcal{A} = 0$ lineage.

\textbf{Part 2.} Suppose $\mathcal{A} > 0$ but $\eta_M = 0$: $(F,K)$ is effectively fixed despite drive. Reduces to Part 1. Additionally, $C_n$ is bounded (Def.~\ref{def:encoding}); without $\eta_M > 0$, no action reaches promotion threshold $\theta_i$ (Theorem~\ref{thm:encoding}), $L_n$ fills, $\rho \to 0$, inference fails. Selection removes $E$.

\textbf{Part 3.} $K(0) = \emptyset$: survives if $\mathcal{A} > 0, \eta_M > 0$. $F(0) = \{\text{id}\}$: the identity action returns feedback to bootstrap. $C_n = C_{min}$: survives if graph structure keeps $\rho = 1$. Substrate, $\sigma$: free parameters. $\square$

\section{Related Work}

\textbf{Free Energy Principle.} FEP~[Friston2010] derives uncertainty minimization from self-organization. This framework does not contradict FEP; FEP agents instantiate it with $U$ as variational free energy and $M$ as Bayesian updating. The framework adds: (i) why $\mathcal{A} > 0$ is structurally necessary (Theorem~\ref{thm:necessity}); (ii) structural constraints on $K$ from survival pressure (Theorem~\ref{thm:retention}); (iii) formal degradation of top-down precision allocation under bounded context (Theorem~\ref{thm:topdown}).

\textbf{Formal theories of intelligence.} AIXI~[Hutter2005] assumes unbounded computation; the central constraint here is bounded $C_n$. Legg and Hutter~[LeggHutter2007] define intelligence as performance; here it is rate of improvement ($\eta_M$). Chollet~[Chollet2019] measures skill-acquisition efficiency --- structurally close to $\eta_M$ --- but treats prior knowledge as a confound; here prior $K$ is constitutive.

\textbf{Bounded rationality.} Simon~[Simon1955] established satisficing under constraints; Griffiths et al.~[GriffithsEtAl2015] formalize resource rationality. The retention criterion (Theorem~\ref{thm:retention}) is a satisficing criterion derived from survival pressure. Theorem~\ref{thm:topdown} is the resource-rational argument that the optimal architecture changes under bounded $C_n$.

\textbf{Cognitive architectures.} Global Workspace Theory~[Baars1988, Dehaene2011] is the escalation architecture in cognitive science language; Theorem~\ref{thm:escalation} derives the escalation criterion from $U$-reduction cost. Gibson's affordances~[Gibson1979] are $K$ indexed by $F$; Theorem~\ref{thm:retention} provides the selection criterion Gibson lacks. Clark~[Clark2015] synthesizes predictive processing with embodied cognition --- hierarchical error minimization under bounded context, the cognitive framing of Theorem~\ref{thm:escalation}.

\textbf{Information theory.} Shannon~[Shannon1948] grounds $U = H(w \mid K)$; Cover and Thomas~[Cover2006] ground $G \geq 0$. Tishby~[Tishby2011] treats the retention threshold as a design parameter; here it is derived. Schmidhuber~[Schmidhuber2010] formalizes curiosity as compression progress ($\approx \eta_M$); here $\mathcal{A} > 0$ is a necessary condition, not a design objective.

\section{Discussion and Conclusion}

The formal results establish a minimal structural account of intelligence under bounded context from four primitives: $W$ with positive entropy, bounded $C_n$, $U_{lethal}^d$, and growing $\mathcal{T}_{reachable}$. From these: $K$ is graph-structured (Theorem~\ref{thm:graph}); escalation dominates top-down allocation (Theorem~\ref{thm:topdown}); the encoding hierarchy follows from cost structure (Theorem~\ref{thm:encoding}); and $\mathcal{A} > 0, \eta_M > 0$ are the unique structural invariants (Theorem~\ref{thm:necessity}).

The framework is substrate-independent. The biological architecture is one instance: heritable encoding at $L_0$, individual learning at $L_n$, cultural transmission as external $K$. The theory predicts that graph structure emerges under bounded-context optimization (confirmed in the engineering system); that top-down orchestration degrades under scaling; and that encoding permanence scales with failure cost. An internal-model account predicts decelerating knowledge growth; civilizational knowledge has accelerated --- consistent with externalized graph-structured $K$ and escalation at civilizational scale.

\textbf{Limitations.} $\mathcal{A}$ and $\eta_M$ are properties of the entity-world coupling measurable by an external observer, not quantities the entity computes in real time. The theory's predictions follow from structural constraints regardless. Extensions to collective intelligence, superintelligence, and sentience are developed in a companion paper; the origin of $\mathcal{A}$ remains open.

\bigskip
\noindent\textbf{Extended version.} Full proofs, additional corollaries, and open questions: [URL].

\end{document}
